\documentclass[12pt]{article}

% =======================================================
% 1. ENGINE + TYPOGRAPHY
%
% Font presets are controlled by the 'fontset' YAML variable.
% Valid values: humanities | demography | methods | (unset = default)
%
%   fontset: humanities  -> EB Garamond         (XeLaTeX/LuaLaTeX)
%                           mathpazo + warning  (pdfLaTeX)
%   fontset: demography  -> XITS                (XeLaTeX/LuaLaTeX)
%                           TeX Gyre Termes     (pdfLaTeX)
%   fontset: methods     -> Source Serif 4 + Fira Code (XeLaTeX/LuaLaTeX)
%                           mathpazo + Fira Code or LM Mono (pdfLaTeX)
%   (default/unset)      -> TeX Gyre Pagella    (XeLaTeX/LuaLaTeX)
%                           mathpazo            (pdfLaTeX)
%
% XeLaTeX/LuaLaTeX fonts are loaded from a local fonts/ directory so the
% bundle is self-contained regardless of system font installation.
% The 'fontpath' YAML variable overrides the default 'fonts/' relative path
% with an absolute path --- required when using the managed template layout
% (e.g., template: ~/Templates/jnc-tex/jnctemplate.tex).
% Example: fontpath: /Users/yourname/Templates/jnc-tex/fonts/
% =======================================================
\usepackage{iftex}

%% --- FONT PRESET MACRO (set from YAML, used for \ifx branching below) ---
\def\jncfontset{demography}
\makeatletter
\def\jncfontset@h{humanities}
\def\jncfontset@d{demography}
\def\jncfontset@m{methods}
%% Boolean flags — set once here, usable everywhere without \makeatletter
\newif\ifjnchumanities
\newif\ifjncdemography
\newif\ifjncmethods
\ifx\jncfontset\jncfontset@h\jnchumanitiestrue\else\jnchumanitiesfalse\fi
\ifx\jncfontset\jncfontset@d\jncdemographytrue\else\jncdemographyfalse\fi
\ifx\jncfontset\jncfontset@m\jncmethodstrue\else\jncmethodsfalse\fi

%% --- FONT PATH (override with 'fontpath:' YAML variable for managed templates) ---
\def\jncfontsdir{fonts/}

%% --- FONT LOADING ---

\ifPDFTeX
  %% pdfLaTeX path: uses TeX distribution packages only (no font files needed).
  \usepackage[T1]{fontenc}
  \ifjncdemography
    %% demography: TeX Gyre Termes (Times-family); ships with every TeXLive/MiKTeX.
    \usepackage{tgtermes}
  \else
    \usepackage{mathpazo}
    \ifjnchumanities
      %% humanities: EB Garamond requires XeLaTeX/LuaLaTeX; mathpazo used instead.
      \PackageWarning{jnctemplate}{fontset=humanities: EB Garamond requires
        XeLaTeX/LuaLaTeX. Falling back to mathpazo under pdfLaTeX.}
    \fi
    \ifjncmethods
      %% methods: try Fira Code for monospace; fall back to Latin Modern Mono.
      \IfFileExists{FiraCode.sty}{%
        \usepackage{FiraCode}%
      }{%
        \PackageWarning{jnctemplate}{fontset=methods: Fira Code not found in the
          TeX distribution. Using Latin Modern Mono. Install the 'firacode'
          package (TeXLive: tlmgr install firacode).}%
      }
    \fi
  \fi
\else
  %% XeLaTeX / LuaLaTeX path: uses fontspec with bundled fonts from fonts/ dir.
  %% Each preset uses \IfFileExists to fail gracefully if fonts/ is not populated.
  %% Run scripts/get-fonts.sh to download font files into fonts/.
  \usepackage{fontspec}
  \ifjnchumanities
    \IfFileExists{\jncfontsdir EBGaramond-Regular.otf}{%
      \setmainfont{EBGaramond}[
        Path           = \jncfontsdir,
        UprightFont    = *-Regular,
        BoldFont       = *-Bold,
        ItalicFont     = *-Italic,
        BoldItalicFont = *-BoldItalic,
        Extension      = .otf]%
    }{%
      \PackageWarning{jnctemplate}{fontset=humanities: \jncfontsdir
        EBGaramond-Regular.otf not found. Run scripts/get-fonts.sh, or set
        fontpath: in YAML if using a managed template. Falling back to default.}%
      \IfFontExistsTF{TeX Gyre Pagella}{%
        \setmainfont{TeX Gyre Pagella}%
      }{%
        \PackageWarning{jnctemplate}{TeX Gyre Pagella not found; using Latin Modern Roman.}%
        \setmainfont{Latin Modern Roman}%
      }%
    }
  \else\ifjncdemography
    \IfFileExists{\jncfontsdir XITS-Regular.otf}{%
      \setmainfont{XITS}[
        Path           = \jncfontsdir,
        UprightFont    = *-Regular,
        BoldFont       = *-Bold,
        ItalicFont     = *-Italic,
        BoldItalicFont = *-BoldItalic,
        Extension      = .otf]%
    }{%
      \PackageWarning{jnctemplate}{fontset=demography: \jncfontsdir
        XITS-Regular.otf not found. Run scripts/get-fonts.sh, or set
        fontpath: in YAML if using a managed template. Falling back to default.}%
      \IfFontExistsTF{TeX Gyre Pagella}{%
        \setmainfont{TeX Gyre Pagella}%
      }{%
        \PackageWarning{jnctemplate}{TeX Gyre Pagella not found; using Latin Modern Roman.}%
        \setmainfont{Latin Modern Roman}%
      }%
    }
  \else\ifjncmethods
    \IfFileExists{\jncfontsdir SourceSerif4-Regular.otf}{%
      \setmainfont{SourceSerif4}[
        Path           = \jncfontsdir,
        UprightFont    = *-Regular,
        BoldFont       = *-Bold,
        ItalicFont     = *-It,
        BoldItalicFont = *-BoldIt,
        Extension      = .otf]%
    }{%
      \PackageWarning{jnctemplate}{fontset=methods: \jncfontsdir
        SourceSerif4-Regular.otf not found. Run scripts/get-fonts.sh, or set
        fontpath: in YAML if using a managed template. Falling back to default.}%
      \IfFontExistsTF{TeX Gyre Pagella}{%
        \setmainfont{TeX Gyre Pagella}%
      }{%
        \PackageWarning{jnctemplate}{TeX Gyre Pagella not found; using Latin Modern Roman.}%
        \setmainfont{Latin Modern Roman}%
      }%
    }
    \IfFileExists{\jncfontsdir FiraCode-Regular.ttf}{%
      \setmonofont{FiraCode}[
        Path        = \jncfontsdir,
        UprightFont = *-Regular,
        BoldFont    = *-Bold,
        Extension   = .ttf]%
    }{%
      \PackageWarning{jnctemplate}{fontset=methods: \jncfontsdir
        FiraCode-Regular.ttf not found. Run scripts/get-fonts.sh to install
        Fira Code. Using default monospace font.}%
    }
  \else
    %% Default or unrecognized fontset: TeX Gyre Pagella with graceful fallback.
    \ifx\jncfontset\empty\else
      \PackageWarning{jnctemplate}{Unknown fontset value '\jncfontset'.
        Valid values: humanities, demography, methods. Using default.}
    \fi
    \IfFontExistsTF{TeX Gyre Pagella}{%
      \setmainfont{TeX Gyre Pagella}%
    }{%
      %% ISSUE 1 FIX: emit warning when TeX Gyre Pagella is unavailable.
      \PackageWarning{jnctemplate}{TeX Gyre Pagella not found on this system.
        Using Latin Modern Roman as fallback. Install tex-gyre for the intended
        Palatino-family appearance.}%
      \setmainfont{Latin Modern Roman}%
    }
  \fi\fi\fi
\fi
\makeatother

\usepackage{microtype}
\usepackage{amsmath,amssymb}
\providecommand{\citeproc}[2]{#2}

% =======================================================
% 1b. COLOR PALETTE
% Loaded early so geometry, titlesec, and hyperref can all use jncaccent.
% methods:    deep tech-blue accent with lighter rule tone
% default/others: black (no visual change)
% =======================================================
\usepackage{xcolor}
\ifjncmethods
  \definecolor{jncaccent}{HTML}{1B4F72}
  \definecolor{jncrule}{HTML}{2E86C1}
\else
  \definecolor{jncaccent}{HTML}{000000}
  \definecolor{jncrule}{HTML}{000000}
\fi

% =======================================================
% 2. PAGE LAYOUT
% humanities:  wider, book-like margins
% methods:     tighter, dense technical report margins
% default/demography: standard 1-inch
% =======================================================
\ifjnchumanities
  \usepackage[left=1.3in,right=1.3in,top=1.1in,bottom=1.1in]{geometry}
\else\ifjncmethods
  \usepackage[left=0.9in,right=0.9in,top=0.9in,bottom=0.9in]{geometry}
\else
  \usepackage[left=1in,right=1in,top=1in,bottom=1in]{geometry}
\fi\fi
\usepackage{setspace}

  %% humanities: slightly more generous leading; methods: tighter
  \ifjnchumanities
    \setstretch{1.55}
  \else\ifjncmethods
    \setstretch{1.25}
  \else
    \onehalfspacing
  \fi\fi

%% humanities: wider paragraph indent for book-like feel
\ifjnchumanities
  \setlength{\parindent}{18pt}
\else
  \setlength{\parindent}{15pt}
\fi
\setlength{\parskip}{0pt}

% =======================================================
% 3. GRAPHICS + FLOATS
% =======================================================
\usepackage{graphicx}
\usepackage{float}
\usepackage{longtable,booktabs}
\usepackage{array}
\usepackage{caption}

\captionsetup{
  font=small,
  labelfont=bf
}

\setcounter{topnumber}{2}
\setcounter{bottomnumber}{1}
\setcounter{totalnumber}{3}

\renewcommand{\arraystretch}{1.1}

% More deterministic float spacing
\setlength{\intextsep}{0.75\baselineskip}
\setlength{\textfloatsep}{1.25\baselineskip}

% --- Auto-scale images ---
\makeatletter
\def\maxwidth{\ifdim\Gin@nat@width>\linewidth\linewidth\else\Gin@nat@width\fi}
\def\maxheight{\ifdim\Gin@nat@height>\textheight\textheight\else\Gin@nat@height\fi}
\makeatother
\setkeys{Gin}{width=\maxwidth,height=\maxheight,keepaspectratio}

% =======================================================
% 4. TABLE AUTOMATION (Pandoc longtable)
% =======================================================
\usepackage{etoolbox}

\AtBeginEnvironment{longtable}{
  \footnotesize
  \singlespacing
  \setlength{\LTleft}{0pt}
  \setlength{\LTright}{0pt}
  \setlength{\LTcapwidth}{\textwidth}
  %% Treat special characters as ordinary text inside table cells.
  %% Pandoc does not escape these when converting markdown tables,
  %% so variable names (underscores), exponents (carets), and
  %% percentage values (percent signs) cause LaTeX errors or silent
  %% data loss. Catcode 12 (other) makes them literal characters
  %% without affecting math mode or comments outside this environment.
  \catcode95=12    %% underscore  (_): subscript operator
  \catcode94=12    %% caret      (^): superscript operator
  \catcode37=12    %% percent    (%): comment char -- silent data loss if unescaped
}

% =======================================================
% 4b. SPECIAL CHARACTER HANDLING IN BODY TEXT
% =======================================================
%% \usepackage{underscore} makes _ safe in ordinary text mode
%% (e.g. variable names mentioned inline in prose) while
%% preserving its subscript role inside math environments.
%% This complements the longtable catcode fix above.
\usepackage{underscore}

% =======================================================
% 5. SECTION FORMATTING + NUMBERING TOGGLE
%
% default:    bold, no rule
% humanities: small-caps, generous spacing
% demography: bold + thin rule below (journal style)
% methods:    bold + accent color + rule below (futuristic)
% =======================================================
\usepackage{titlesec}

%% Named rule commands for \titleformat after-code.
%% Necessary because \titleformat's optional [after-code] argument uses [],
%% and LaTeX's optional-arg scanner doesn't handle nested [] depth, so
%% \titleformat{...}{}{}{}{}[\titlerule[0.5pt]] would close the outer []
%% at the first ] it encounters. Wrapping in a command avoids nesting.
\newcommand{\jncdemorule}{\titlerule[0.3pt]}
\newcommand{\jncmetrule}{\color{jncrule}\titlerule[0.5pt]}

  \setcounter{secnumdepth}{0}
  \setcounter{tocdepth}{0}
  \ifjnchumanities
    \titleformat{\section}{\large\scshape}{}{0pt}{}
    \titleformat{\subsection}{\normalsize\bfseries\itshape}{}{0pt}{}
    \titleformat{\subsubsection}{\normalsize\itshape}{}{0pt}{}
    \titlespacing{\section}{0pt}{2.5ex plus 0.5ex}{1ex}
  \else\ifjncdemography
    \titleformat{\section}{\large\bfseries}{}{0pt}{}[\jncdemorule]
    \titleformat{\subsection}{\normalsize\bfseries}{}{0pt}{}
    \titleformat{\subsubsection}{\normalsize\bfseries\itshape}{}{0pt}{}
  \else\ifjncmethods
    \titleformat{\section}{\large\bfseries\color{jncaccent}}{}{0pt}{}[\jncmetrule]
    \titleformat{\subsection}{\normalsize\bfseries\color{jncaccent}}{}{0pt}{}
    \titleformat{\subsubsection}{\normalsize\itshape\color{jncaccent}}{}{0pt}{}
    \titlespacing{\section}{0pt}{1.75ex plus 0.5ex}{0.5ex}
  \else
    \titleformat{\section}{\large\bfseries}{}{0pt}{}
    \titleformat{\subsection}{\normalsize\bfseries}{}{0pt}{}
    \titleformat{\subsubsection}{\normalsize\bfseries\itshape}{}{0pt}{}
  \fi\fi\fi

% =======================================================
% 6. HYPERREF + HEADERS/FOOTERS
% =======================================================
%% methods: colored links reinforce the accent palette
\ifjncmethods
  \usepackage[colorlinks=true,linkcolor=jncaccent,urlcolor=jncrule,citecolor=jncaccent]{hyperref}
\else
  \usepackage[hidelinks]{hyperref}
\fi
\usepackage{fancyhdr}

% =======================================================
% 7. OPTIONAL MODES
% =======================================================
\newif\ifjncanon
\jncanonfalse

\newcommand{\jncmaincolumns}{1}


% =======================================================
% 8. PANDOC SUPPORT
% =======================================================
\providecommand{\tightlist}{%
  \setlength{\itemsep}{0pt}%
  \setlength{\parskip}{0pt}%
}

\providecommand{\pandocbounded}[1]{#1}


%% NOTE (Issue 5): knitr/Quarto may inject booktabs, float, longtable, or
%% other packages already loaded above via header-includes. This is benign ---
%% LaTeX's \usepackage deduplicates package loading --- but may produce harmless
%% "already loaded" notes in verbose log output. No action required.
\usepackage{booktabs}
\usepackage{pdflscape}

% --- Pandoc table width arithmetic support ---
\usepackage{calc}
\providecommand{\real}[1]{#1}

% =======================================================
% RUNNING TITLE + PAGESTYLE
% =======================================================
\pagestyle{fancy}
\setlength{\headheight}{14.5pt}
\fancyhf{}
\fancyfoot[C]{\thepage}

% Ensure "plain" pages (e.g., first page after \section in some classes) match
\fancypagestyle{plain}{
  \fancyhf{}
  \fancyfoot[C]{\thepage}
}

\fancyhead[R]{\small\itshape LastName: Short Running Title}

% =======================================================
% 9. CSL REFERENCES (Pandoc)
% =======================================================
\newlength{\cslhangindent}
\setlength{\cslhangindent}{1.5em}

\providecommand{\citeproctext}{}

\newenvironment{CSLReferences}[2]
 {%
  %% ISSUE 3 FIX: revert to single column before references when maincolumns: 2.
  %% \twocolumn persists through the body into the bibliography section otherwise.
  %% \onecolumn forces a page break (standard behavior for two-column papers).
  \ifnum\jncmaincolumns=2\relax\onecolumn\fi
  \section*{References}
  \begin{list}{}{%
   \setlength{\itemindent}{0pt}
   \setlength{\leftmargin}{0pt}
   \setlength{\parsep}{0pt}
   \ifodd #1
    \setlength{\leftmargin}{\cslhangindent}
    \setlength{\itemindent}{-\cslhangindent}
   \fi
   \ifnum #2 > 0
    \setlength{\itemsep}{#2\baselineskip}
   \else
    \setlength{\itemsep}{0pt}
   \fi
  }%
  \renewcommand{\bibitem}[2][]{\item}
 }
 {\end{list}}

% =======================================================
% 10. R / SOFTWARE MACROS
% =======================================================
\newcommand{\pkg}[1]{\texttt{#1}}
\newcommand{\CRANpkg}[1]{\texttt{#1}}
\newcommand{\code}[1]{\texttt{#1}}

% =======================================================
% BEGIN DOCUMENT
% =======================================================
\begin{document}

% =============================================================
% TITLE PAGE
% methods:  left-aligned, thick accent rule top and bottom
% others:   centered (humanities/demography/default)
% =============================================================

%% Shared author-block macro.
%%
%% Author names come from the for(author) loop — Quarto passes author.name
%% through its normalize-author.lua filter. All other author sub-fields
%% (department, institution, email, orcid, etc.) are stripped by that filter
%% and never reach the Pandoc template.
%%
%% Solution: supply contact details via top-level author_* YAML variables,
%% which Quarto does not process and therefore passes through unchanged.
%%
%%   author_dept:   Department of Sociology
%%   author_inst:   Queens College, City University of New York
%%   author_addr:   65-30 Kissena Blvd., Queens, NY 11367, USA
%%   author_email:  joseph.cohen@qc.cuny.edu
%%   author_web:    https://jncohen.commons.gc.cuny.edu
%%   author_orcid:  0000-0002-6197-4453
%%
\newcommand{\jncauthorblock}{%
\ifjncanon\else
\textbf{Roni Pitkowsky}\par
{\small\setlength{\parskip}{1pt}%
Department of X\par%
CUNY Queens College\par%
Street Address, City, State ZIP\par%
\texttt{\href{mailto:you@university.edu}{\nolinkurl{you@university.edu}}}\par%
\url{https://yourwebsite.edu}\par%
ORCID: 0000-0000-0000-0000\par%
}%
\fi
}

\ifjncmethods
%% -------------------------------------------------------
%% METHODS: futuristic left-aligned title page
%% -------------------------------------------------------
\begin{titlepage}
\thispagestyle{empty}
\noindent\textcolor{jncaccent}{\rule{\textwidth}{4pt}}\par
\vspace{3em}
\begin{flushleft}
\begin{singlespace}

{\setstretch{0.9}\LARGE\bfseries My Paper Title\par}

\vspace{0.5em}
{\large\color{jncaccent} Optional subtitle\par}

\vspace{2em}

\jncauthorblock

\ifjncanon\else
\vspace{0.25em}
{\small\color{jncaccent}\textit{April 25, 2026}}\par
\fi

\end{singlespace}
\end{flushleft}

\vfill

\noindent\textcolor{jncrule}{\rule{0.35\textwidth}{0.5pt}}\par
\vspace{0.5em}
\begin{minipage}{0.85\textwidth}
\small\singlespacing
\textbf{Abstract}\par\vspace{0.25em}
Your abstract here (150-250 words recommended).
\end{minipage}
\vspace{1em}

\ifjncanon\else
\begin{minipage}{0.85\textwidth}
\footnotesize\singlespacing
\noindent\textcolor{jncaccent}{\rule{0.25\textwidth}{0.4pt}}\par
\textbf{Acknowledgements}\par
Optional. Funding, data sources, etc.
\end{minipage}
\vspace{1em}
\fi

\noindent\textcolor{jncaccent}{\rule{\textwidth}{2pt}}\par
\end{titlepage}

\else
%% -------------------------------------------------------
%% DEFAULT / HUMANITIES / DEMOGRAPHY: centered title page
%% -------------------------------------------------------
\begin{titlepage}
\thispagestyle{empty}

\vspace*{0.06\textheight}
\begin{center}

{\setstretch{0.9}\LARGE\bfseries My Paper Title\par}

\vspace{0.25em}
{\large Optional subtitle\par}

\vspace{2em}

\begin{singlespace}
{\normalsize

\jncauthorblock

% --- DATE ---
%% NOTE (Issue 4): To preserve a human-readable date string, use inline R:
%%   date: "`r format(Sys.Date(), '%B %d, %Y')`"
\ifjncanon\else
\vspace{0.5em}
\textit{April 25, 2026}\par
\fi

}
\end{singlespace}

\vfill

\begin{minipage}{0.8\textwidth}
\small\singlespacing
\textbf{Abstract}\par\vspace{0.25em}
Your abstract here (150-250 words recommended).
\end{minipage}

\vfill

\ifjncanon\else
\begin{minipage}{0.8\textwidth}
\footnotesize\singlespacing
\noindent\rule{0.3\textwidth}{0.4pt}\par
\textbf{Acknowledgements}\par
Optional. Funding, data sources, etc.
\end{minipage}
\fi

\end{center}
\end{titlepage}
\fi

\clearpage
\setcounter{page}{1}

% =======================================================
% MAIN BODY
% =======================================================

% Main text column layout toggle
\ifnum\jncmaincolumns=2\relax
  \twocolumn
\fi

\ifjncanon
\begingroup
\small
\noindent\textit{This manuscript is set to anonymous mode. Identifying metadata fields are suppressed on the title page.}\par
\endgroup
\vspace{1em}
\fi

\section{Methods}\label{methods}

\subsection{Data and Analytic
Approach}\label{data-and-analytic-approach}

Data for this analysis comes from the U.S. Federal Reserve's (2022)
\emph{Survey of Consumer Finances} (SCF) (Board of Governors of the
Federal Reserve System, 2022). The SCF is a triennial survey of U.S.
households, which collects data on household assets, debts, incomes and
expenditures, as well as other information. The survey utilizes a
complex sampling mechanism with 999 replicate weights and a multiple
imputation scheme (five implicates) to handle missing data and estimate
variance. To facilitate the handling of this data, I utilized the
\texttt{scf} R package (Cohen, 2025).

The unit of analysis is the household. This report examines \textbf{high
debt burden}, defined as a Payment-to-Income Ratio (PIR) exceeding 40\%.
Categorical predictors (wealth, education, age, sex) were recoded to
facilitate interpretation. Key predictors include age group, sex,
marital status, wealth group, education, presence of children,
homeownership, and race/ethnicity.

\subsection{Theoretical Framework: Testing a Structural
Hypothesis}\label{theoretical-framework-testing-a-structural-hypothesis}

This analysis is guided by emerging scholarship on the structural
production of financial vulnerability, particularly Frederick Wherry's
concept of ``viral debt'' (Gardner, Gray, \& Wherry, 2026). Wherry
argues that indebtedness circulates through communities along pathways
shaped by historical exclusion and policy design, rather than spreading
randomly or reflecting individual choices alone. From this perspective,
racial disparities in debt burden should persist even after accounting
for measurable socioeconomic differences---because financial strain is
produced by structural barriers, not just compositional factors.

To test this hypothesis, I estimated multivariate regression models that
isolate the independent association of race with high debt burden after
controlling for wealth, education, family structure, and homeownership.
The key question is: Do racial disparities persist when comparing
households with similar economic profiles? If yes, this would suggest
that debt burden operates through racialized structures rather than
individual characteristics.

\phantomsection\label{refs}
\begin{CSLReferences}{1}{0}
\bibitem[\citeproctext]{ref-FederalReserve2022}
Board of Governors of the Federal Reserve System. (2022). \emph{Survey
of consumer finances (SCF)} {[}Data set{]}. Retrieved from
\url{https://www.federalreserve.gov/econres/scfindex.htm}

\bibitem[\citeproctext]{ref-Cohen2025}
Cohen, J. N. (2025). \emph{Scf: Tools for analyzing the survey of
consumer finances}. Retrieved from \url{https://github.com/jncohen/scf}

\bibitem[\citeproctext]{ref-Wherry2026}
Gardner, J., Gray, M., \& Wherry, F. F. (Eds.). (2026). \emph{Viral
debt: The production and reproduction of economic vulnerability}. New
York: Routledge. \url{https://doi.org/10.4324/9781003481461}

\end{CSLReferences}

% =======================================================
% REFERENCES
% =======================================================
\clearpage

\end{document}